% Cover letter for the initial submission to
% International Journal for Numerical Methods in Engineering (I{}JNME, Wiley).
% Build: pdflatex cover-letter
% Fill the [TO FILL] fields (editor name, date, ORCID line) before submitting.

\documentclass[11pt]{letter}
\usepackage[a4paper,margin=2.5cm]{geometry}
\usepackage[T1]{fontenc}
\usepackage[utf8]{inputenc}
\usepackage{newtxtext}
\usepackage{xcolor}
\usepackage[colorlinks=true,linkcolor=blue!55!black,urlcolor=blue!55!black]{hyperref}

\signature{Mathias Rodr\'iguez Castro}
\address{Facultad de Ingenier\'ia\\Universidad de la Rep\'ublica\\Montevideo, Uruguay\\\href{mailto:mathiasr@fing.edu.uy}{mathiasr@fing.edu.uy}}

\begin{document}

\begin{letter}{The Editors\\International Journal for Numerical Methods in Engineering}

\opening{Dear Editors,}

We submit for your consideration the manuscript \emph{``Coverage versus Selectivity in
Pre-Export Row Scaling for Generated Mixed-Integer Programs''} as a Research Article.

Model generators build mixed-integer programs (MIPs) component by component and know,
before export, which rows are local to a component, which variables it owns, and which
rows couple components---structure that the flat coefficient matrix handed to the solver
discards. The paper asks whether this pre-export metadata can be reused as numerical
information to drive positive row scaling, a transformation that changes the numbers the
solver sees while leaving the feasible set, objective, and integer restrictions exactly
unchanged. We believe the work fits I{}JNME's scope in numerical methods for engineering
computation: it concerns the numerical representation of generated engineering
optimization models and the interaction between that representation and commercial
solvers, evaluated on a hydrothermal-dispatch platform used as a demanding
block-structured testbed rather than as a domain application.

The central contribution is methodological as much as algorithmic. Across three instance
classes, two commercial solvers, and two optimality tolerances, we separate four kinds of
evidence that are easily conflated---algebraic equivalence, numerical representation,
solver-facing conditioning, and matched per-solver runtime---and show that they need not
agree. Positive row scaling is exact; structure-aware scaling robustly improves
solver-facing numerical diagnostics; and runtime improves under Gurobi but not robustly
under CPLEX. A role-blind control then isolates the runtime lever as a cheap per-row
normalization rather than the row-role metadata, with a class-dependent outcome that we
trace to a coverage gap over the highest-magnitude rows. The result is a clean, controlled
demonstration that the rule which best conditions a matrix is not the one that solves
fastest---a cautionary and, we hope, useful lesson for numerical preprocessing of
generated MIPs. We report the negative and mixed results (the CPLEX neutrality, and the
class in which the structure-blind control loses) as fully as the positive ones.

The work is original and is not under consideration elsewhere; no part has been previously
published. The manuscript has a single author, who is responsible for all of its content.
A reproducibility package---cleaned source code, the redistributable dispatch and synthetic
instances, aggregate per-scenario results, and the scripts that rebuild every table and
figure---is publicly available and archived with a DOI, as stated in the Data Availability
Statement; the experiments were carried out on the ClusterUY high-performance computing
platform. The manuscript declares no conflicts of interest and received no specific
funding.

We would be glad to suggest suitable reviewers on request, and we thank you for
considering our submission.

\closing{Sincerely,}

\end{letter}
\end{document}
